\documentclass[../mathNotesPreamble]{subfiles}

\providecommand{\relscalefact}{1.4}
\begin{document}
\relscale{\relscalefact}
  \section{8.5: Partial Order Relations}
  \begin{defn*}
    Let $R$ be a relation on a set $A$. $R$ is \textbf{antisymmetric} if, and only if, for every $a$ and $b$ in $A$, if $\rel{a}{b}$ and $\rel{b}{a}$, then $a=b$.
    \vspace*{\baselineskip}

    Alternatively, if $a\neq b$, then either $\notrel{a}{b}$ or $\notrel{b}{a}$.
  \end{defn*}
  \begin{ex*}
    Let $R_1$ and $R_2$ be relations on $\set{0,1,2}$ defined as follows:
      \[R_1=\set{(0,2),(1,2),(2,0)}, \qquad R_2=\set{(0,0),(0,1),(0,2),(1,1),(1,2)}\]
    Draw the directed graphs for $R_1$ and $R_2$ and indicate which relations are antisymmetric.
    \vspace*{\stretch{1}}
  \end{ex*}
  \pagebreak

  \begin{defn*}
    Let $R$ be a relation defined on a set $A$. $R$ is a \textbf{partial order relation} if, and only if, $R$ is reflexive, antisymmetric, and transitive.
  \end{defn*}
  \begin{ex*}
    For non-empty sets $U$, and $V$, define the relation $S$ such that
      \[\rel[S]{U}{V} \Leftrightarrow U\subseteq V\]
    Prove that $S$ is a partial order relation.
    \vspace*{\stretch{1}}
  \end{ex*}
  \begin{ex*}
    For integers $a$, and $b$, define the relation $D$ such that
      \[\rel[D]{a}{b} \Leftrightarrow a\mid b\]
    Prove that $D$ is a partial order relation.
    \vspace*{\stretch{1}}
  \end{ex*}
  \pagebreak
  \begin{defn*}
    The symbol $\preceq$ is often used to denote a general partial order relation.
  \end{defn*}

  \begin{thmBox*}[Lexicographic Order]
    Let $A$ be a set with a partial order relation $R$, and let $S$ be a set of strings over $A$. Define a relation $\preceq$ on $S$ as follows:
    \begin{quote}
      Let $s$ and $t$ be any strings in $S$ of length $m$ and $n$, respectively, where $m$ and $n $ are positive integers, and let $s_m$ and $t_m$ be the characters in the $m$th position for $s$ and $t$, respectively.
    \end{quote}
    \begin{enumerate}
      \item If $m\leq n$ and the first $m$ characters of $s$ and $t$ are the same, then $s\preceq t$.
      \item If the first $k-1$ characters in $s$ and $t$ are the same, $\rel{s_k}{t_k}$, and $s_k\neq t_k$, then $s\preceq t$.
      \item If $\lambda$ is the null string, then $\lambda \preceq s$.
    \end{enumerate}
    If no strings are related by $\preceq$ other than by these three conditions, then $\preceq$ is a partial order relation on $S$.

    This partial order relation is called the \textbf{lexicographic order for $S$} that corresponds to the partial order $R$ on $A$.
  \end{thmBox*}
  \pagebreak

  \begin{ex*}
    Let $A=\set{x,y}$ and let $R$ be the following partial order relation on $A$:
      \[R=\set{(x,x),(x,y),(y,y)}\]
    \begin{extasks}[after-item-skip=\stretch{1}](2)
      \task* Is $x\preceq x$? \hspace*{\stretch{1}}
        Is $x\preceq xx$? \hspace*{\stretch{1}}
        Is $yx\preceq yxy$? \hspace*{\stretch{1}}
      \task Is $xxxyyy\preceq xy$?
      \task Is $x\preceq y$?
      \task Is $\lambda \preceq xy$?
      \task Is $xyy\preceq xyx$?
    \end{extasks}
    \vspace*{\stretch{1}}
  \end{ex*}
  \pagebreak

  \begin{defn*}
    A \textbf{Hasse diagram} is a graphical representation for a partial relation. To obtain a Hasse diagram, start with a directed graph of the relation with vertices placed so that all arrows point upwards. Then eliminate
    \begin{itemize}
      \item the loops at all vertices,
      \item all arrows implied by the transitive property, and
      \item the direction indicators on the arrows.
    \end{itemize}
  \end{defn*}
  \begin{ex*}
    Let $A=\set{1,2,3,9,18}$. Consider the partial order relation ``divides'':
      \[\rel[D]{a}{b} \Leftrightarrow a\mid b\]

    \newcommand{\parRelGraph}{%
      \vspace*{1.5\baselineskip}
      \newline \centering
      \begin{tikzpicture}[
        every node/.style={circle, fill=black, inner sep=1.5pt},
        label distance=2pt
      ]
        \node[label=-45:1] at (0,0) {};
        \node[label=-135:2] at (-2,2) {};
        \node[label=-45:3] at (2,2) {};
        \node[label=0:9] at (3.5,3.5) {};
        \node[label=180:18] at (0,5) {};
      \end{tikzpicture}
    }
    \noindent
    \begin{minipage}{0.475\linewidth}
      Draw the directed graph of $D$:
      \parRelGraph
    \end{minipage}\hspace*{\stretch{1}}%
    \begin{minipage}{0.475\linewidth}
      Generate the Hasse diagram of $D$:
      \parRelGraph
    \end{minipage}
  \end{ex*}
  \pagebreak

  \begin{ex*}
    Consider the ``subset'' relation, $\subseteq$, on the set $\mathscr{P}\parens{\set{a,b,c}}$:
      \[\rel{U}{V} \Leftrightarrow U\subseteq V\]
    Construct the Hasse diagram for this relation.
    \vspace*{\stretch{1}}
  \end{ex*}
  \pagebreak

  \begin{defn*}
    Suppose $\preceq$ is a partial order relation on a set $A$. Elements $a$ and $b$ of $A$ are said to be \textbf{comparable} if, and only if, either $a\preceq b$ or $b\preceq a$. Otherwise, $a$ and $b$ are called \textbf{noncomparable}. \vspace*{0.75\baselineskip}

    When all the elements of a partial order relation are comparable, the relation is called a \textbf{total order}. \vspace*{0.75\baselineskip}

    If $R$ is a partial order relation on a set $A$, and for any two elements $a$ and $b$ in $A$ either $\rel{a}{b}$ or $\rel{b}{a}$, then $R$ is a \textbf{total order relation} on $A$. \vspace*{0.75\baselineskip}

    A set $A$ is called a \textbf{partially ordered set} (or \textbf{poset}) with respect to a relation $\preceq$ if, and only if, $\preceq$ is a partial order relation on $A$. \vspace*{0.75\baselineskip}

    A set $A$ is called a \textbf{totally ordered set} with respect to a relation $\preceq$ if, and only if $A$ is partially ordered with respect to $\preceq$ and $\preceq$ is a total order. \vspace*{0.75\baselineskip}

    A subset $B$ of $A$ is called a \textbf{chain} if, and only if, the elements in each pair of elements in $B$ are comparable. The \textbf{length of a chain} is one less than the number of elements in the chain. \vspace*{0.75\baselineskip}

    \emph{Note}: The Hasse diagram for a total order relation can be drawn as a single vertical chain.
  \end{defn*}
  \begin{ex*}
    The set $\mathscr{P}(\set{a,b,c})$ is partially ordered with respect to the subset relation. Find a chain of length $3$ in $\mathscr{P}(\set{a,b,c})$.
    \vspace*{\stretch{1}}
  \end{ex*}
  \pagebreak

  \begin{defn*}
    Let a set $A$ be partially ordered with respect to a relation $\preceq$. An element in $A$ is called a
    \begin{enumerate}
      \item \textbf{maximal element of $A$} $\Leftrightarrow \forall b\in A$, either $b\preceq a$ or $b$ and $a$ are not comparable.
      \item \textbf{greatest element of $A$} $\Leftrightarrow \forall b\in A$, $b\preceq a$.
      \item \textbf{minimal element of $A$} $\Leftrightarrow \forall b\in A$, either $a\preceq b$ or $b$ and $a$ are not comparable.
      \item \textbf{least element of $A$} $\Leftrightarrow \forall b\in A$, $a\preceq b$.
    \end{enumerate}
  \end{defn*}
  \begin{ex*}
    Let $A=\set{a,b,c,d,e,f,g,h,i}$ have the partial ordering $\preceq$ defined by the following Hasse diagram. Find all maximal, minimal, greatest, and least elements of $A$.
    \begin{center}
      \begin{tikzpicture}[node distance=15mm and 20mm,
        every node/.style={circle, fill=black, inner sep=1.5pt},
        label distance=2pt
      ]
        \node[label=180:a] (a) at (0,0) {};
        \node[below=of a, label=180:b] (b) {};
        \node[below left = of b, label=-30:c] (c) {};
        \node[below right = of b, label=0:d] (d) {};
        \node[above right = of d, label=0:e] (e) {};
        \node[above = of e, label=0:f] (f) {};
        \node[above = of f, label=0:g] (g) {};
        \node[above right = of e, label=0:h] (h) {};
        \node[below = of h, label=0:i] (i) {};

        \draw (a) -- (b)  -- (d) -- (e) -- (f) -- (g) -- (a);
        \draw (b) -- (c) (e) -- (h) -- (g) (h) -- (i);
      \end{tikzpicture}
    \end{center}
  \end{ex*}

  \pagebreak
\end{document}
